\chapter{Width comparison and finite extinction}
\label{ch:extinction-v4}

We now turn the deformation of loop families into a finite extinction
argument. A nonzero third homotopy class gives a nonnegative width.
Ordinary Ricci flow forces this width below a scalar profile, and surgery
cannot cause an upward jump from the left. One profile therefore bounds
the width through finitely many events. Its eventual negativity contradicts
the existence of a surviving component.

The initial class and its width must be fixed before the terminal time and
surviving component are chosen. We use the simply connected, time-zero
branch: the initial manifold already has trivial second homotopy and
nonzero third homotopy. The later-start argument in
\cite[Section 18.3.2, pp. 431--432]{MT2007} is not needed here.

\section{Widths and the deformation input}

On an ordinary slab, $M$ is a compact, connected, Hausdorff, second-countable
smooth three-manifold without boundary, $x\in M$, and $\pi_2(M,x)=0$.
Write $\Lambda^1M$ for the $C^1$ free-loop space of
Chapter~\ref{ch:filling-width-v4}.

\begin{definition}[Null family]
\label{def:v4-null-family}
\lean{PoincareMT.M61NullFamily}
A null family is a continuous map $F:S^2\to\Lambda^1M$ whose every member
is nullhomotopic. No continuously chosen family of spanning disks is required.
\source{MT2007}\dcref{Definition 18.17, p. 430}
\end{definition}

For a null loop $\gamma$, let $A_g(\gamma)$ be the infimum of the areas of
its Lipschitz spanning disks, with boundary equal to $\gamma$ up to the
admissible boundary reparameterization. Section~\ref{sec:filling-short-loops}
supplies such disks, nonnegative filling area, and continuity on null loops.

\begin{definition}[Family and class width]
\label{def:v4-family-width}
\uses{def:v4-null-family}
\label{def:v4-free-width}
\lean{PoincareMT.m61FamilyWidth}
\lean{PoincareMT.m61FreeClassWidthRange}
\lean{PoincareMT.m61FreeClassWidth}
For a null family $F$, set
\[
 W_g(F)=\max_{c\in S^2}A_g(F(c)),\qquad
 W_g[F]=\inf\{W_g(G):G\text{ is null and freely homotopic to }F\}.
\]
The maximum exists for each fixed family; the class infimum need not be
attained. Proposition~\ref{prop:filling-loop-identification} labels this
same competitor class by $\alpha\in\pi_2(\Lambda^1M,\mathrm{const}_x)$
and identifies $\alpha$ with $\xi\in\pi_3(M,x)$. We also write
$W_g(\alpha)$ for this free-class infimum.
\source{MT2007}\dcref{Claim 18.16 and Definition 18.17, p. 430}
\end{definition}

\begin{proposition}[Maximum and near minimizers]
\label{prop:v4-width-properties}
\uses{prop:filling-width-properties}
\uses{def:v4-family-width}
\lean{PoincareMT.m61FamilyWidth_from_M60}
\lean{PoincareMT.m61Widths}
These widths are finite and nonnegative. For every $e>0$ there is a null
family $G$ freely homotopic to $F$ with $W_g(G)<W_g[F]+e$.
\end{proposition}
\begin{proof}
The continuous function $c\mapsto A_g(F(c))$ attains its maximum on $S^2$.
The competitor range contains $F$ and is bounded below by zero. If no
competitor had width less than $W_g[F]+e$, that larger number would also
be a lower bound, contradicting the definition of the infimum.
\end{proof}

Two separate short-loop inputs will be used. For a fixed metric there is
$r>0$ such that a null family all of whose members have length less than
$r$ is freely homotopic to a constant family. For every $\eta>0$ there
is $q(\eta)>0$ such that a loop shorter than $q(\eta)$ bounds a disk of
area less than $\eta$. In the first assertion, contracting each loop to
its initial point produces a sphere of constant loops; $\pi_2(M)=0$
then contracts that sphere of points. These are the results of
Section~\ref{sec:filling-short-loops}.

For a Ricci flow $g(t)$ on $[a,b]$, define
\[
 B_a(t)=\int_a^t R_{\min}(u)/2\,du,\qquad
 \Phi_{a,A}(t)=e^{-B_a(t)}
       \left(A-2\pi\int_a^t e^{B_a(u)}\,du\right).
\]
The scalar-curvature infimum $R_{\min}$ on the compact slice is continuous.
Consequently
\[
 \Phi_{a,A}(a)=A,\quad
 \partial_t\Phi_{a,A}=-2\pi-\tfrac12R_{\min}\Phi_{a,A},\quad
 \Phi_{a,A+d}(b)=\Phi_{a,A}(b)+e^{-B_a(b)}d.
\]
Theorem~\ref{thm:loop-family-deformation} gives, for each null family $G$
and $z>0$, null families $D_t$ freely homotopic to $G$ such that
\[
 |A_{g(a)}(D_a(c))-A_{g(a)}(G(c))|<z,
\]
and for \emph{each} $c$ either
\[
 \operatorname{Length}_{g(b)}(D_b(c))<z
 \quad\hbox{or}\quad
 A_{g(b)}(D_b(c))\leq\Phi_{a,A_{g(a)}(D_a(c))}(b)+z.
\]
Different members may satisfy different alternatives. The disk first
variation and Gauss--Bonnet argument producing $-2\pi$ is part of that
deformation theorem.
\cite[Definition 18.23, Proposition 18.24 and Claim 18.26, pp. 433--434]{MT2007}

\section{Comparison on an ordinary slab}

\begin{proposition}[Endpoint comparison]
\label{prop:v4-endpoint-comparison}
\uses{prop:v4-width-properties,thm:loop-family-deformation,prop:filling-short-area,prop:filling-loop-identification}
\lean{PoincareMT.m66_deformed_comparison}
\lean{PoincareMT.m66_endpoint_comparison}
If $F$ represents a nonzero class and $a\leq b$, then
\[
 W_{g(b)}[F]\leq\Phi_{a,W_{g(a)}[F]}(b).
\]
\end{proposition}
\begin{proof}
Write $w=W_{g(a)}[F]$, $A=\Phi_{a,w}(b)$ and $E=e^{-B_a(b)}>0$.
Fix $e,\eta>0$ and choose an $e$-near minimizing family $G$.
Choose the terminal short-family radius $r$ and small-disk radius
$q(\eta)$, and only then set $z=\min\{e,r,q(\eta)\}$ and deform $G$.
For every member the initial error gives
$A_{g(a)}(D_a(c))\leq w+e+z$. Positivity and affinity of the profile
in its initial value show that each member in the area alternative has
\[
 A_{g(b)}(D_b(c))\leq A+E(e+z)+z\leq A+(2E+1)e.
\]
At least one member is not short: otherwise the short-family theorem
would make $D_b$, and hence $F$, trivial. Its nonnegative area gives
$0\leq A+(2E+1)e$. This does not exclude short members.
Every short member has filling area less than $\eta$ by the small-disk
estimate. Taking the family maximum and then the class infimum gives
\[
 W_{g(b)}[F]\leq\max\{\eta,A+(2E+1)e\}.
\]
Letting $e\downarrow0$ in the nonnegativity inequality gives $A\geq0$.
Given $\epsilon>0$, set $\eta=\epsilon$ and $e=\epsilon/(2E+1)$.
The last display is at most $A+\epsilon$. Let $\epsilon\downarrow0$.
\end{proof}

Continuity requires an estimate uniform over all disks and families.

\begin{lemma}[Metric comparison for width]
\label{lem:v4-metric-width}
\uses{prop:v4-width-properties}
\lean{PoincareMT.m66_freeClassWidth_le_of_metric_le}
If metrics on the compact manifold satisfy $h(v,v)\leq Cg(v,v)$
for $C>0$, then $W_h[F]\leq CW_g[F]$.
\end{lemma}
\begin{proof}
For tangent vectors $u,v$ with $u\ne0$, replace $v$ by
$v-g(u,v)u/g(u,u)$. Each Gram determinant is unchanged, and the
vectors become $g$-orthogonal. Bound the two diagonal $h$ terms by
$C$ times the $g$ terms and discard the nonpositive cross-square:
\[
 \det\operatorname{Gram}_h(u,v)
 \leq C^2\det\operatorname{Gram}_g(u,v).
\]
The case $u=0$ is immediate. Square roots give the area-density factor
$C$. The identity map is $\sqrt C$-Lipschitz, so the same disk remains
admissible, with unchanged boundary and almost-everywhere differential.
Its new area is at most $C$ times its old area. Compactness supplies a
reverse metric bound too, so admissibility is equivalent for both metrics.
Taking disk infima, family maxima and class infima proves the result;
the competitor class is unchanged.
\end{proof}

\begin{proposition}[Smooth-time estimates]
\label{prop:v4-smooth-time}
\uses{prop:v4-endpoint-comparison,lem:v4-metric-width}
\lean{PoincareMT.m66_smooth_time_theory}
\lean{PoincareMT.m66Width_continuousAt}
\lean{PoincareMT.m66_forward_difference_bound}
For the nonzero class above, $w(t)=W_{g(t)}[F]$ is nonnegative and
continuous on $[a,b]$. For every $s\in[a,b)$ and $\epsilon>0$ there
is $\delta>0$ such that, if $s<t\leq b$ and $t<s+\delta$, then
\begin{equation}
 \frac{w(t)-w(s)}{t-s}
 \leq-2\pi-\frac{R_{\min}(s)}2w(s)+\epsilon.
 \label{eq:v4-forward}
\end{equation}
\source{MT2007}\dcref{Proposition 18.18, p. 431}
\dcref{Claim 18.21, pp. 432--433}
\end{proposition}
\begin{proof}
Curvature continuity on the compact space-time slab bounds
$|\operatorname{Ric}(v,v)|\leq Dg(v,v)$ for some $D\geq0$.
Integrating $\partial_tg=-2\operatorname{Ric}$ on a fixed tangent
vector gives $g(t)\leq e^{2D|t-s|}g(s)$ and the reverse bound. Thus
Lemma~\ref{lem:v4-metric-width} gives
\[
 e^{-2D|t-s|}w(s)\leq w(t)\leq e^{2D|t-s|}w(s).
\]
The squeeze theorem proves continuity, including at slab endpoints.
No minimizing representative is chosen uniformly in time.

Apply Proposition~\ref{prop:v4-endpoint-comparison} to the restricted
slab $[s,t]$. It gives $w(t)\leq\Phi_{s,w(s)}(t)$ with equality of
initial values at $s$. The profile has right derivative
$-2\pi-R_{\min}(s)w(s)/2$ there. Subtract $w(s)$, divide by $t-s>0$
and use the definition of this derivative to obtain
\eqref{eq:v4-forward}. Differentiability of $w$ is not asserted.
\end{proof}

\section{One class through a surgery history}

Fix a component path $X_t\subset M_t$ on $[0,T]$. Its components are
compact connected smooth three-manifolds with basepoints $x_t$; at an
event the slice means the outgoing slice. There are finitely many events,
none at zero. On $(a,b]$ without events, actual slice diffeomorphisms
$r_{a,b}:X_a\to X_b$ have compatible composition and basepoint transport.
At an event $s$, an earlier preterminal time $s^-<s$ has no event in
$(s^-,s)$. Theorem~\ref{thm:comparison-transport} gives a comparison
homotopy equivalence from that parent to the selected child. Its parent
and child are simply connected; we retain this specialization here.

\begin{proposition}[A coherent nonzero class along the history]
\label{prop:v4-class-transport}
\uses{prop:filling-ancestry-transport,prop:filling-loop-identification}
\lean{PoincareMT.m67_finite_transport_section,PoincareMT.m67_class_path_exists}
Fix one based/free identification throughout, $\pi_2(X_0,x_0)=0$,
and $\xi_0\ne0$ in $\pi_3(X_0,x_0)$. Transport ordinarily by the slice
diffeomorphisms and specified basepoint paths. At an event, transport to
the preterminal parent, then by the actual comparison map, and finally
along a path in the child to $x_s$. Homotopy equivalences preserve
triviality of $\pi_2$ and nontriviality of the third homotopy element;
basepoint change is an isomorphism too.
This constructs one compatible nonzero class at every time on the path.
\end{proposition}

\begin{proof}
For completeness, order the finite set consisting of zero and the events.
Define the class at zero. The preterminal time $s^-<s$ of each event
lies after a previously defined anchor, whose class can be transported
ordinarily to $s^-$. Apply the event map to define the new anchor class.
At an arbitrary $t$, transport from the largest anchor not exceeding
$t$. Two times with no intervening event have the same last anchor, so
the composition rule proves compatibility of their classes. The recursive
definition gives the event rule. No unrelated class is chosen on a child.
\end{proof}

Let $\alpha_t$ be the resulting loop-space class and put
$w(t)=W_{g(t)|X_t}(\alpha_t)$. The based/free identification makes this
the actual free-class width on each ordinary slab. An interval's initial
width is evaluated at that interval's start, not silently replaced by
the time-zero width.

\begin{lemma}[The event class is independent of smoothing tolerance]
\label{lem:v4-event-class}
\uses{prop:v4-class-transport,thm:comparison-transport}
Area estimates use smooth approximants of an event map. All these
approximants have the same based $\pi_3$ map at the fixed comparison
basepoint. Naturality of $\pi_2(\Lambda^1X)\simeq\pi_3(X)$ shows that
their postcomposed families give the same $\alpha_s$ after rebasing
along the fixed child path. Changing the approximation tolerance therefore
does not change the width being estimated.
\lean{PoincareMT.m67_event_class_datum_exists}
\end{lemma}

\section{The estimate at an event}

The strict event bounds are $\delta(s)<\delta_{\rm cut}$ and
$h(s)<R_{\rm cut}^{-1/2}$ for the cutoffs of the local surgery theorem;
the flow accuracy must also satisfy its terminal-neighborhood threshold.
Section~\ref{sec:global-comparison-calibration} makes these choices before
constructing the flow. Once these bounds hold, the comparison
theorem gives the following uniform conclusion. For each $\eta>0$
there is $\delta_0>0$ such that every preterminal $t$ with
$s-\delta_0<t<s$ admits a smooth $(1+\eta)$-Lipschitz map
$f_t:X_t\to X_s$ inducing the same based homotopy map as above.
These are the same-map and same-class conclusions of
Theorem~\ref{thm:comparison-transport}.

\begin{lemma}[Disk postcomposition]
\label{lem:v4-disk-transport}
\uses{lem:v4-metric-width}
\lean{PoincareMT.m67_filling_transport_of_lipschitz}
If $f:(X,g)\to(Y,h)$ is smooth and $L$-Lipschitz, $L>0$, then
$A_h(f\circ\gamma)\leq L^2 A_g(\gamma)$ for every null loop $\gamma$.
\end{lemma}
\begin{proof}
Its differential satisfies $h(df(v),df(v))\leq L^2g(v,v)$.
The Gram-determinant argument of Lemma~\ref{lem:v4-metric-width},
applied to these differentials, bounds the area Jacobian by $L^2$.
For an admissible disk $D$, the disk $f\circ D$ is Lipschitz, retains
the boundary reparameterization, and is differentiable almost everywhere
by the chain rule. Its area density is integrable and its area is at most
$L^2\operatorname{Area}_g(D)$. Take the infimum over $D$.
\end{proof}

Postcomposition sends each pre-event competitor for $\alpha_t$ to a
competitor for $\alpha_s$. Taking family maxima and then pre-event
infima in the disk estimate yields
\begin{equation}
 w(s)\leq(1+\eta)^2w(t)\quad(s-\delta<t<s),\qquad
 \delta=\min\{\delta_0,s-s^-\}>0.
 \label{eq:v4-event-factor}
\end{equation}
The minimum puts every test time on the actual preterminal regular
piece. The tolerance is chosen before the neighborhood.

\begin{proposition}[Changing width estimates]
\label{prop:v4-event-width}
\uses{prop:v4-smooth-time,lem:v4-disk-transport,lem:v4-event-class,thm:comparison-transport,prop:filling-loop-identification}
\lean{PoincareMT.m67_surgery_width_theory}
\lean{PoincareMT.m67_event_factor_bounds}
\lean{PoincareMT.m67_lower_limit_of_factor_bounds}
The constructed width is nonnegative, obeys \eqref{eq:v4-forward} on
ordinary pieces, is continuous at regular times, and is right-continuous
at events. For every event $s$ and $\epsilon>0$ there is $\delta>0$
such that
\[
 s-\delta<t<s\quad\Longrightarrow\quad w(s)\leq w(t)+\epsilon.
\]
This is the meaning of $w(s)\leq\liminf_{t\uparrow s}w(t)$ here.
No incoming limit is assumed to exist.
\source{MT2007}\dcref{Proposition 18.18, p. 431}
\dcref{Claim 18.22, p. 433}
\end{proposition}
\begin{proof}
Nonnegativity and regular estimates apply to the same widths on ordinary
pieces. Finiteness of the event set gives an outgoing regular interval at
each time before $T$. Slab continuity, including its left endpoint,
proves right continuity. At $T$ the assertion within $[0,T]$ is immediate.

Put $v=w(s)\geq0$, choose
$\eta=\min\{1,\epsilon/(4(v+1))\}>0$, and then choose the neighborhood
in \eqref{eq:v4-event-factor}. If $w(t)\geq v$ there is nothing to prove.
Otherwise $0\leq w(t)<v$, and
\[
 v-w(t)\leq(2\eta+\eta^2)w(t)\leq3\eta v
             <4\eta(v+1)\leq\epsilon.
\]
This holds uniformly for all pre-times in that neighborhood.
\end{proof}

\section{Scalar comparison without differentiating the width}

The normalized scalar-curvature bound is
\begin{equation}
 R_{\min}(t)\geq-\frac6{1+4t}.
 \label{eq:v4-scalar-clock}
\end{equation}
Since $w\geq0$, \eqref{eq:v4-forward} bounds its upper forward slopes
by $F(t,w(t))=-2\pi+3w(t)/(1+4t)$. Nonnegativity is essential in
this substitution.

\begin{proposition}[Regular-slab forward comparison]
\label{prop:v4-m81}
\uses{prop:rf-evolution}
\lean{PoincareMT.m81RegularSlabComparison}
Let $a\leq b$, with $1+4t>0$ on $[a,b]$. Suppose $f,G$ are
continuous there, $f(a)\leq G(a)$, $G'=F(t,G)$, and for every
$s\in[a,b)$ and $\epsilon>0$, all sufficiently near $t>s$ satisfy
$(f(t)-f(s))/(t-s)\leq F(s,f(s))+\epsilon$.
Then $f\leq G$ throughout the slab.
\end{proposition}
\begin{proof}
Put $k=3/(1+4a)+1>0$. Since the clock increases,
$3/(1+4s)<2k$. For $\eta>0$ use the barrier
$H(t)=G(t)+\eta e^{2k(t-a)}$. At a contact $f(s)=H(s)$,
\[
 F(s,f(s))-H'(s)
 =\left(\frac3{1+4s}-2k\right)\eta e^{2k(s-a)}<0.
\]
The forward quotient bound and the right derivative of $H$ imply
$f(t)-H(t)<0$ for every sufficiently near $t>s$.

If $f(t_*)>H(t_*)$, continuity and $f(a)<H(a)$ give a zero of $f-H$
before $t_*$. The zero set in $[a,t_*]$ is compact, so it has a largest
element $s<t_*$. Continuity and absence of further zeros make $f-H$
positive on $(s,t_*]$, contrary to the negative values immediately to
the right of $s$. Thus $f\leq H$. The perturbation is bounded by
$\eta e^{2k(b-a)}$, so letting $\eta\downarrow0$ proves $f\leq G$.
Only right slopes at a last contact were used, not a derivative of $f$.
\end{proof}

\begin{lemma}[The scalar comparison profile]
\label{lem:v4-profile}
\lean{PoincareMT.m68_profile_equation}
For $T_1\geq0$, the scalar solution with value $w_1$ there is
\begin{equation}
 P_{T_1,w_1}(t)=
 w_1\left(\frac{1+4t}{1+4T_1}\right)^{3/4}
 +2\pi(1+4T_1)^{1/4}(1+4t)^{3/4}-2\pi(1+4t).
 \label{eq:v4-profile}
\end{equation}
It satisfies $P'= -2\pi+3P/(1+4t)$ and $P(T_1)=w_1$ for $t\geq T_1$.
\end{lemma}
\begin{proof}
Multiplying the ODE by $(1+4t)^{-3/4}$ gives derivative
$-2\pi(1+4t)^{-3/4}$, whose integral is
$-2\pi((1+4t)^{1/4}-(1+4T_1)^{1/4})$. Solving for $P$ gives
\eqref{eq:v4-profile} and $P(T_1)=w_1$. All real powers have positive
arguments.
\end{proof}

\section{Restarting through finitely many events}

\begin{lemma}[Restart at an event]
\label{lem:v4-event-restart}
\uses{prop:v4-event-width,lem:v4-profile}
\lean{PoincareMT.m68_event_restart_bound}
Fix $[T_1,T_2]\subset[0,T]$ and retain the single profile
$P=P_{T_1,w(T_1)}$ on that whole interval. Suppose $w(t)\leq P(t)$
for $a<t<s$ and $s$ is an event. Then $w(s)\leq P(s)$.
\end{lemma}
\begin{proof}
For $\epsilon>0$, use the event
inequality with error $\epsilon/2$ and continuity of $P$ with error
$\epsilon/2$. Choose $t\in(a,s)$ in both neighborhoods. Then
\[
 w(s)\leq w(t)+\epsilon/2\leq P(t)+\epsilon/2<P(s)+\epsilon.
\]
Letting $\epsilon\downarrow0$ gives $w(s)\leq P(s)$ without any
incoming width limit.
\end{proof}

\begin{proposition}[Finite-event scalar comparison]
\label{prop:v4-scalar-clock}
\uses{prop:v4-m81,prop:v4-event-width,lem:v4-event-restart,lem:v4-profile}
\lean{PoincareMT.m68_scalar_clock_statement}
For every $t\in[T_1,T_2]$, $0\leq w(t)\leq P_{T_1,w(T_1)}(t)$.
\end{proposition}
\begin{proof}
If $[a,b]$ has no event in $(a,b)$ and $w(a)\leq P(a)$, then
ordinary comparison gives the bound at $b$ when $b$ is regular.
When $b$ is an event, compare on each $[a,c]$, $a<c<b$, and then
use the preceding restart argument. Outgoing continuity at $a$ allows
$a$ itself to be an event.

Induct on a finite set covering all interior events. The empty case is
the preceding argument. Remove an element $s$. If $s\notin(a,b)$,
the remaining set still covers all interior events. If $a<s<b$, apply
the induction hypothesis to $[a,s]$ and then $[s,b]$. Each interval
has its interior events covered by the smaller set, and the first bound
supplies $w(s)\leq P(s)$ for the second. Starting from equality at
$T_1$ proves the result on $[T_1,t]$ for every $t$. The profile and
its initial width are never reset.
\end{proof}

\begin{proposition}[Finite-piece propagation and contradiction]
\label{prop:v4-finite-piece}
\label{prop:v4-component-contradiction}
\uses{prop:v4-scalar-clock}
\lean{PoincareMT.m69FinitePiecePropagation}
\lean{PoincareMT.m70FiniteExtinctionContradiction}
The transported classes remain nonzero and their actual filling widths
satisfy the preceding bound. A component path therefore cannot exist
through $B\in[T_1,T_2]$ with $P_{T_1,w(T_1)}(B)<0$.
\end{proposition}
\begin{proof}
The class construction preserves nontriviality and identifies the starting
width with the free infimum at $T_1$. Scalar comparison would give
$0\leq w(B)\leq P_{T_1,w(T_1)}(B)<0$. This contradicts existence
of that path; a common extinction time still requires common initial data.
\end{proof}

\section{A common negative time and the first empty event}

Fix the global flow of Section~\ref{sec:global-assembly}, with time domain
$[0,\infty)$, nonempty initial slice, locally finite events, permanence of
emptiness, and no surgery strictly after an empty slice. We require the
following ancestry data from Section~\ref{sec:filling-fixed-initial-class}.
There is one fixed initial component, basepoint, induced metric and
nonzero class $\alpha_0$, chosen before the observation time. Every
target slice is covered by the terminal components of finitely many paths
starting at this same initial component. Each path has the comparison
and class data above. All scalar and strict event bounds hold for this
same flow.

\begin{lemma}[One initial width for all observed paths]
\label{lem:v4-common-width}
\uses{prop:v4-class-transport,prop:filling-component-ancestry}
The class construction is anchored at $\alpha_0$. Hence every such
path, for every terminal time, has
\[
 w(0)=w_0:=W_{g(0)|X_0}(\alpha_0).
\]
Equality of initial components also identifies their basepoints and
induced metrics; an unspecified isomorphism of components would not
suffice for this equality of widths.
\lean{PoincareMT.m71WidthChoice_initial_width}
\end{lemma}

\begin{lemma}[A common negative comparison time]
\label{lem:v4-negative-time}
\uses{lem:v4-common-width,lem:v4-profile,prop:v4-width-properties}
Set
\begin{equation}
 c=2+\frac{w_0}{2\pi},\qquad T_*=(c^4-1)/4.
 \label{eq:v4-extinction-time}
\end{equation}
Since $w_0\geq0$, $c\geq2$ and $T_*>0$. The zero-start profile is
$(w_0+2\pi)(1+4t)^{3/4}-2\pi(1+4t)$, and therefore
\[
 P_{0,w_0}(T_*)=(w_0+2\pi)c^3-2\pi c^4=-2\pi c^3<0.
\]
Here $w_0=2\pi(c-2)$ and $c>0$ justify the equality.
\lean{PoincareMT.m71ProfileAlgebra_neg}
\end{lemma}

\begin{theorem}[Global finite extinction]
\label{prop:v4-global-extinction}
\uses{thm:global-controlled-flow,prop:filling-component-ancestry,prop:filling-ancestry-transport,prop:foundations-initial,prop:foundations-projective-exclusion}
\uses{prop:v4-finite-piece,lem:v4-negative-time}
\lean{PoincareMT.m71GlobalFiniteExtinction}
\lean{PoincareMT.m71FirstEmptySurgery}
Under the global-flow and fixed-initial-class ancestry hypotheses just
stated, the slice at $T_*$ is empty. An actual surgery time $S\leq T_*$
has empty outgoing slice, every earlier slice is nonempty, and all later
slices are empty.
\source{MT2007}\dcref{Theorem 18.1, p. 415}
\dcref{Section 18.3.2, pp. 431--432}
\end{theorem}
\begin{proof}
A surviving point at $T_*$ belongs, by the target cover, to a component
path reaching $T_*$. Its initial width is $w_0$, so the common negative
profile contradicts Proposition~\ref{prop:v4-component-contradiction}.
Thus the slice at $T_*$ is empty.

Take the maximum $S$ of zero and all surgery times in $[0,T_*]$.
This is a nonempty finite set, with no surgery in $(S,T_*]$.
If $S<T_*$, ordinary-slab identification carries emptiness back to
$S$; if $S=T_*$ it is already empty. Initial nonemptiness excludes
$S=0$, so $S$ is an actual event. An empty slice at an earlier $s<S$
would prohibit that event by the no-surgery-after-emptiness property.
Thus all earlier slices are nonempty, and permanence gives later emptiness.
\end{proof}

The last event before an already empty observation thus locates the first
empty slice; no attainment of an infimum of empty times was assumed.
This event is the input to finite-history reconstruction. The uniformity
here comes from one time-zero class on the simply connected initial
manifold. The broader argument in \cite[pp. 431--432]{MT2007} starts after
second homotopy has vanished and takes a bound over finitely many initial
components at that later time. That additional route is not asserted here.
